% year: תשפ"ג
% type: מבחן
% semester: א
% moed: ב
% number: 4ב
% points: 10
% categories: taylor
% source: exams/מבחנים/תשפג/מועד ב תשפג.docx

\begin{question}
השתמשו בפולינום טיילור מסדר שני כדי להעריך את הביטוי $\sqrt[4]{82}$ והעריכו את השגיאה.
\end{question}

\begin{hint}
פתחו את $f(x)=x^{1/4}$ סביב $x_0=81$ והשתמשו בשארית לגרנז'.
\end{hint}

\begin{solution}
$f(x)=x^{1/4}$, $x_0=81$, $f(81)=3$.
\[ f'(x)=\tfrac14x^{-3/4},\quad f''(x)=-\tfrac{3}{16}x^{-7/4},\quad f'''(x)=\tfrac{21}{64}x^{-11/4}. \]
\[ f'(81)=\frac{1}{4\cdot27}=\frac{1}{108},\qquad f''(81)=-\frac{3}{16\cdot 3^7}=-\frac{1}{16\cdot729}=-\frac{1}{11664}. \]
פולינום טיילור מסדר 2 סביב $81$, בנקודה $x=82$ ($x-x_0=1$):
\[ \sqrt[4]{82}\approx P_2(82)=3+\frac{1}{108}-\frac{1}{2\cdot11664}=3+\frac{1}{108}-\frac{1}{23328}=\frac{70199}{23328}\approx 3.0092164. \]
\textbf{הערכת שגיאה} (שארית לגרנז'): קיימת $c\in(81,82)$ כך ש-
\[ R_2=\frac{f'''(c)}{3!}(82-81)^3=\frac{21}{64\cdot 6}\,c^{-11/4}. \]
$c^{-11/4}$ יורדת, ולכן $c^{-11/4}<81^{-11/4}=3^{-11}$, ומכאן
\[ 0<R_2<\frac{21}{384\cdot 3^{11}}=\frac{21}{384\cdot177147}\approx 3.1\cdot10^{-7}. \]
כלומר $\sqrt[4]{82}\approx3.0092164$ בשגיאה קטנה מ-$3.1\cdot 10^{-7}$ (והקירוב קטן מעט מהערך האמיתי, $3.0092167\ldots$).
\end{solution}
